\section{Descritivo do projeto}
\label{sec:descritivo}

\subsection{Contexto}

A Sociedade Brasileira de Tecnologia e Computação (SBTC) promove congressos, simpósios, workshops e escolas regionais em todo o país. Por muitos anos, inscrições, pagamentos, credenciamento e certificados ficaram espalhados em sistemas e planilhas diferentes. Com o aumento do número de eventos e de participantes, esse arranjo ficou difícil de manter, e a SBTC decidiu centralizar tudo em um sistema próprio, o EventOS.

O caso traz situações que tornam o banco de dados menos trivial do que parece. O valor de uma inscrição depende da categoria, do lote, da modalidade e da situação de associação do participante na data. Algumas categorias exigem documento comprobatório antes do pagamento. Atividades têm vagas limitadas, instituições associadas têm benefícios com limites, trabalhos científicos precisam de um autor inscrito e certificados só podem ser emitidos para quem cumpriu os critérios da edição. E tudo precisa ficar rastreável, porque o sistema guarda dados pessoais e financeiros.

Este capítulo descreve a primeira entrega do projeto (Marco 1): a delimitação do minimundo, a modelagem conceitual e lógica e a implantação do banco no PostgreSQL. O código e os scripts estão em um repositório público, disponível em
\repo.

\subsection{Objetivos}

O objetivo geral é construir um banco de dados relacional que sustente o ciclo de vida completo de uma inscrição em evento científico da SBTC, do cadastro do participante à emissão do certificado, mantendo o histórico financeiro e a rastreabilidade das operações.

Os objetivos específicos acompanham os processos do caso:
\begin{enumerate}
  \item representar eventos, edições, eventos co-localizados e a oferta de cada edição (categorias, faixas de preço e atividades);
  \item calcular e congelar o valor devido por inscrição, separando inscrição, associação, atividades, taxas de publicação e descontos;
  \item controlar a validação de documentos comprobatórios antes de autorizar o pagamento;
  \item registrar tentativas de pagamento, cancelamentos e reembolsos segundo a política de cada edição;
  \item controlar a capacidade das atividades, inclusive quando várias pessoas disputam a última vaga;
  \item aplicar e limitar os benefícios de instituições associadas;
  \item verificar a cobertura de trabalhos científicos por inscrições de autores;
  \item registrar credenciamento e presença e emitir certificados apenas para participantes aptos;
  \item guardar a versão de cada termo aceito, a auditoria das operações relevantes e regras de acesso a dados pessoais.
\end{enumerate}

\subsection{Extensão proposta}

O enunciado pede que a equipe estenda o caso com requisitos próprios. A Tabela~\ref{tab:re} lista os 16 requisitos estendidos (RE) que ficaram no projeto. Os de RE1 a RE10 surgiram da leitura crítica do caso, ainda na primeira versão do modelo. Os de RE11 a RE17 apareceram durante a revisão do escopo, quando o grupo percebeu lacunas que o modelo reduzido precisava cobrir. O RE7 foi retirado no corte de escopo e não entra na contagem; mantivemos a numeração original para que cada requisito possa ser rastreado desde a primeira versão, e a justificativa da retirada está na Tabela~\ref{tab:simplificados}. A maior parte dos requisitos foi proposta em discussão do grupo; a coluna Proposto por indica quando a proposta partiu de um integrante.

\begin{longtable}{@{}l L{4.9cm} L{4.8cm} L{2.8cm}@{}}
\caption{Requisitos estendidos propostos pelo grupo}\label{tab:re}\\
\toprule
RE & Requisito & Como o banco atende & Proposto por\\
\midrule
\endfirsthead
\toprule
RE & Requisito & Como o banco atende & Proposto por\\
\midrule
\endhead
\bottomrule
\endfoot
RE1 & O valor cobrado fica congelado e aponta para o preço de origem & item financeiro com \cod{vl_item} e inscrição ligada ao preço (Enquadra) & Grupo\\
RE2 & O preço fica reservado por um prazo e expira depois dele & \cod{nr_dias_reserva} na edição, \cod{dt_validade_preco} na inscrição e situação EXPIRADA & Grupo\\
RE3 & Uma inscrição libera as edições vinculadas a ela & derivado do autorrelacionamento Agrega (simplificado) & Grupo\\
RE4 & Reserva de vaga em atividade com prazo e lista de espera & \cod{dt_expira_reserva} e situação RESERVADA; lista de espera retirada (simplificado) & Grupo\\
RE5 & Troca de categoria com recálculo e histórico & novo preço, itens antigos marcados como SUBSTITUIDO e registro na auditoria & Grupo\\
RE6 & Limites de benefício por escopo e controle de uso individual & limites por período e por edição; escopo por evento retirado (simplificado) & Grupo\\
RE8 & Anonimização sem perda de dados financeiros & \cod{ind_anonimizada} e procedimento de anonimização & Filipe Santana\\
RE9 & Cobertura de trabalho explícita e política de publicação & relacionamento Cobre, política na edição e liberação da cobertura no cancelamento & Grupo\\
RE10 & Perfis de acesso por papel & papéis, permissões por coluna e visão mascarada & Filipe Santana\\
RE11 & Inscrição não paga expira e libera a vaga & situação EXPIRADA e procedimento de expiração & Grupo\\
RE12 & Motivo da rejeição de documento registrado & \cod{ds_motivo_rejeicao} & Grupo\\
RE13 & Distribuição dos participantes por UF & \cod{uf_pessoa} & Matheus França\\
RE14 & Tipo de vínculo com a instituição & \cod{tp_vinculo} & Grupo\\
RE15 & Reembolso registrado no próprio pagamento & \cod{vl_reembolso}, \cod{dt_reembolso} e situação ESTORNADO & Filipe Santana\\
RE16 & Certificado revogável e reemitido sem perder o histórico & \cod{status_certificado} (VALIDO ou REVOGADO) & Túlio Melo\\
RE17 & Atividade restrita a algumas categorias & relacionamento Restringe; atividade sem restrição fica aberta a todas & Felipe Schuler\\
\end{longtable}

\subsection{Minimundo}

O EventOS guarda eventos da SBTC e suas edições anuais. Cada edição tem período, local, modalidade e as políticas que valem para ela: prazo de reserva do preço, regra de reembolso, critérios de certificado e limite de trabalhos por inscrição. Uma edição pode agregar outras edições que acontecem junto com ela.

Cada edição oferece categorias de inscrição e define faixas de preço no tempo, como primeiro lote, segundo lote e inscrição tardia. O preço é a combinação de uma categoria, uma faixa e uma situação de associação (não associado, associado, combo com nova associação ou combo com renovação). Quando a modalidade é combo, o preço separa o valor do evento do valor da associação.

Pessoas têm cadastro único, com CPF ou passaporte, e podem ter períodos de associação à SBTC e vínculos com instituições ao longo do tempo. Uma instituição associada pode oferecer benefícios (gratuidade ou desconto), com limites por período e por edição.

Uma inscrição pertence a uma pessoa e se enquadra em um preço. Ela passa pelas situações iniciada, aguardando validação, aguardando pagamento, paga, concluída, cancelada ou expirada. O valor devido é formado por itens financeiros, gerados no momento da cobrança e nunca alterados depois. A inscrição pode ter documentos comprobatórios, tentativas de pagamento, atividades adicionais selecionadas e respostas aos termos vigentes.

Atividades têm capacidade e preço próprios e podem ser restritas a algumas categorias. Trabalhos científicos pertencem a uma edição, têm um ou mais autores e podem ser cobertos pela inscrição de um autor; o trabalho que passa do limite da inscrição gera uma taxa extra. Durante o evento são registrados o credenciamento e a presença em atividades. Certificados são emitidos para pessoas aptas e podem ser revogados e reemitidos. As operações relevantes ficam registradas na auditoria, com o responsável quando houver.

\subsection{Delimitação do escopo}
\label{sec:escopo}

A primeira versão do modelo tinha 34 entidades (Figura~\ref{fig:merv1}). Ela cobria o caso inteiro, mas tratava como entidade muita coisa que era parâmetro, subtipo ou registro duplicado, e desenhava as ligações como atributos de chave estrangeira, sem os losangos de relacionamento que a notação de Chen exige. O professor indicou que o escopo estava grande demais e sugeriu algo entre 10 e 15 entidades, sem rigidez.

O grupo avaliou cada entidade com uma pergunta: ela tem identidade e ciclo de vida próprios no minimundo, ou é uma característica de outra entidade? O que era característica virou atributo; o que era ligação virou relacionamento. O resultado foram 20 entidades. O grupo não desceu até 15 porque as cinco restantes sustentam requisitos do próprio enunciado: sem Preço não há como combinar categoria, faixa e associação (RF6, RF7); sem Item Financeiro não há como identificar a origem de cada valor (RF12); sem Período de Associação não se sabe a situação do participante em uma data (RF4); sem Benefício não há limites por instituição (RF17); sem Auditoria não há rastreabilidade (RF25).

A Tabela~\ref{tab:removidas} mostra o destino das entidades retiradas da primeira versão. A Tabela~\ref{tab:adicionadas} lista as entidades que não estão na especificação inicial do enunciado e a Tabela~\ref{tab:atributos} resume os atributos alterados em relação a ela.

\begin{longtable}{@{}L{3.4cm} L{4.5cm} L{6.1cm}@{}}
\caption{Entidades da primeira versão que deixaram de ser entidades}\label{tab:removidas}\\
\toprule
Entidade na V1 & Destino no modelo final & Justificativa\\
\midrule
\endfirsthead
\toprule
Entidade na V1 & Destino no modelo final & Justificativa\\
\midrule
\endhead
\bottomrule
\endfoot
Boleto, Cartão, Faturamento e Empenho & \cod{tp_pagamento} e atributos opcionais em Pagamento & O RF14 pede formas diferentes de pagamento, e cada subtipo tinha só dois ou três atributos próprios.\\
Reembolso & \cod{vl_reembolso}, \cod{dt_reembolso} e situação ESTORNADO em Pagamento & O reembolso é a devolução de um pagamento específico, uma relação de um para um.\\
Política de Reembolso & \cod{nr_dias_limite_reembolso} e \cod{pc_reembolso} em Edição & O RF15 define a política por edição.\\
Regra de Certificação & \cod{ind_exige_credenciamento} e \cod{pc_frequencia_minima} em Edição & É um parâmetro da edição, sem ciclo de vida próprio.\\
Modalidade de Inscrição & \cod{tp_associacao} em Preço & É uma das dimensões do preço (RF6, RF7).\\
Sessão e Credenciamento & \cod{dt_credenciamento} em Inscrição e \cod{dt_presenca} em Inscrição em Atividade & A atividade já tem início e fim; a sessão é tratada como atividade (RF20, RF21).\\
Histórico de Status & Auditoria & O RF9 e o RF25 usam o mesmo mecanismo de registro.\\
Usuário Administrativo & Pessoa e relacionamento Organiza (\cod{tp_papel}) & Quem administra também é pessoa; a mesma estrutura atende o RF24 e o responsável do RF25.\\
Uso de Benefício & Relacionamento Facilita (\cod{dt_uso_beneficio} em Inscrição) & Uma inscrição usa no máximo um benefício.\\
Cobertura de Trabalho & Relacionamento Cobre & Um trabalho é coberto por no máximo uma inscrição.\\
Vínculo Institucional repetido & Uma única entidade associativa & A entidade aparecia duas vezes no diagrama.\\
\end{longtable}

\begin{longtable}{@{}L{3.6cm} L{10.7cm}@{}}
\caption{Entidades acrescentadas em relação à especificação do enunciado}\label{tab:adicionadas}\\
\toprule
Entidade & Por que existe\\
\midrule
\endfirsthead
\toprule
Entidade & Por que existe\\
\midrule
\endhead
\bottomrule
\endfoot
Preço & O enunciado diz que o preço depende de edição, categoria, faixa e modalidade de associação e pede que isso seja modelado sem redundância. A associativa entre Categoria e Faixa guarda uma linha por combinação.\\
Período de Associação & O RF4 pede saber se a associação estava vigente em uma data. Um campo de situação na pessoa só diria a situação de hoje.\\
Vínculo Institucional & O RF16 pede vínculo com período de validade. A mesma pessoa pode voltar a uma instituição em outro período, por isso o vínculo é uma entidade associativa.\\
Benefício & O RF17 exige limites por benefício, e a instituição pode ter mais de um benefício.\\
Inscrição em Atividade & O enunciado diz que a relação entre inscrição e atividade tem informações próprias (data, situação, cancelamento). A presença em atividade também fica aqui.\\
Item Financeiro & O RF12 pede identificar a origem de cada valor cobrado, e o RF27 pede que valores antigos continuem reproduzíveis.\\
Auditoria & O RF25 pede o objeto alterado, a operação, o instante e o responsável. A mesma tabela guarda o histórico de situações do RF9.\\
\end{longtable}

\begin{longtable}{@{}L{3.2cm} L{11.1cm}@{}}
\caption{Atributos alterados em relação à especificação do enunciado}\label{tab:atributos}\\
\toprule
Tabela & Alteração e motivo\\
\midrule
\endfirsthead
\toprule
Tabela & Alteração e motivo\\
\midrule
\endhead
\bottomrule
\endfoot
Todas & Chaves com \cod{GENERATED ALWAYS AS IDENTITY} e textos longos com \cod{varchar(4000)}, porque \cod{SERIAL} e \cod{text} não fazem parte do SQL padrão (Regra 6).\\
\cod{tbl_pessoa} & Inclui \cod{uf_pessoa}, porque o relatório xi do RF26 pede distribuição por UF, e \cod{ind_anonimizada} (RF30, RE8).\\
\cod{tbl_edicao_evento} & Inclui a edição agregadora (RF2) e as políticas da edição: reserva, reembolso, credenciamento, frequência mínima, limite de trabalhos e taxa por trabalho extra.\\
\cod{tbl_inscricao} & Retira \cod{ce_edicao} e \cod{ce_categoria}, que são obtidos pelo preço, e \cod{vl_total_inscricao}, que é a soma dos itens. O enunciado pede essa avaliação. Inclui o preço, o benefício usado, a validade do preço e o credenciamento.\\
\cod{tbl_atividade} & Inclui \cod{vl_atividade}, porque a atividade pode ter preço próprio.\\
\cod{tbl_documento_comprovacao} & \cod{ce_usuario_analise} passa a ser \cod{ce_pessoa_analise}, porque o analista é uma pessoa. Inclui o caminho do arquivo (o binário fica fora do banco) e o motivo da rejeição (RE12).\\
\cod{tbl_pagamento} & Inclui vencimento (boleto e empenho) e os dados do reembolso (RE15).\\
\cod{tbl_instituicao} & Inclui \cod{ind_associada}, condição para conceder benefícios (RF17).\\
\cod{tbl_trabalho} & Inclui a inscrição que cobre o trabalho (RF19).\\
\cod{tbl_certificado} & Inclui \cod{status_certificado} para revogação e reemissão (RE16).\\
\end{longtable}

Alguns requisitos foram atendidos com simplificação, e um requisito estendido foi retirado. A Tabela~\ref{tab:simplificados} registra cada caso. Os sete requisitos novos (RE11 a RE17) estão na Tabela~\ref{tab:re}; todos nasceram de atributos que o modelo reduzido já precisava ter e que o grupo tornou explícitos como requisito.

\begin{longtable}{@{}l L{3.4cm} L{9.6cm}@{}}
\caption{Requisitos simplificados ou retirados}\label{tab:simplificados}\\
\toprule
Req. & Situação & Justificativa\\
\midrule
\endfirsthead
\toprule
Req. & Situação & Justificativa\\
\midrule
\endhead
\bottomrule
\endfoot
RF14 & simplificado & As formas de pagamento são um tipo com atributos opcionais e restrições por tipo, sem uma entidade para cada forma.\\
RF15 & simplificado & A política de reembolso é uma faixa única por edição (percentual e prazo), sem tabela escalonada por data.\\
RF21 & simplificado & O enunciado fala em várias ocorrências de presença. O modelo guarda uma presença por atividade e trata cada sessão como atividade, o que evita mais uma entidade e ainda permite calcular a frequência.\\
RF22 & simplificado & Os critérios de certificado ficam na edição. Não há critério próprio por atividade; o certificado de atividade exige a presença registrada.\\
RF27 & simplificado & As políticas da edição não têm versões. Para compensar, elas ficam bloqueadas depois da abertura das inscrições, e os valores já cobrados ficam congelados nos itens.\\
RE3 & simplificado & O acesso às edições agregadas é derivado do relacionamento Agrega; não há escolha de eventos por modalidade.\\
RE4 & simplificado & A reserva de vaga tem prazo, mas a lista de espera foi retirada.\\
RE6 & simplificado & Os limites valem por período e por edição; o limite por evento foi retirado.\\
RE7 & retirado & As regras de certificação deixaram de ser entidade versionada e viraram atributos da edição, protegidos contra alteração depois da abertura das inscrições.\\
\end{longtable}

\subsection{Atendimento dos requisitos do enunciado}

A Tabela~\ref{tab:rf} mostra como cada requisito funcional é atendido no modelo e no banco. Os arquivos citados estão na pasta \cod{sql/} do repositório.

\begin{longtable}{@{}l L{6.2cm} L{6.6cm}@{}}
\caption{Atendimento dos requisitos funcionais}\label{tab:rf}\\
\toprule
RF & No modelo & No banco\\
\midrule
\endfirsthead
\toprule
RF & No modelo & No banco\\
\midrule
\endhead
\bottomrule
\endfoot
RF1 & Edição é entidade fraca de Evento (Realiza) & edição única por evento e ano\\
RF2 & Autorrelacionamento Agrega em Edição & agregação de um nível só, sem ciclo; visão \cod{vw_acesso_inscricao}\\
RF3 & Pessoa com cadastro único; Efetua liga pessoa e inscrições & e-mail, CPF e passaporte com país únicos; CPF ou passaporte obrigatório\\
RF4 & Período de Associação, fraca de Pessoa & períodos sem sobreposição; preço coerente com a situação na data; \cod{vw_situacao_associacao}\\
RF5 & Categoria, fraca de Edição & nome único dentro da edição\\
RF6 & Faixa, fraca de Edição, e Preço como associativa & faixas sem sobreposição; faixa vigente na data da inscrição\\
RF7 & \cod{tp_associacao} e \cod{vl_associacao} em Preço; Origina & item de associação separado; período criado ao pagar\\
RF8 & Documento, fraca de Inscrição (Garante), e Analisa & documento nasce pendente, muda uma vez e não é excluído; pagamento só com documento aprovado\\
RF9 & \cod{status_inscricao} com sete situações & transições controladas por trigger; histórico na auditoria\\
RF10 & Inscrição em Atividade e Restringe & atividade da mesma edição ou de edição agregada; categoria permitida\\
RF11 & \cod{nr_capacidade} em Atividade; ocupação derivada & contagem de reservas e confirmações antes de aceitar\\
RF12 & Item Financeiro (Compõe, Tarifa, Sobretaxa) & itens gerados por trigger; \cod{vw_valor_inscricao}\\
RF13 & Pagamento, fraca de Inscrição (Transaciona) & quitação só com pagamentos confirmados; \cod{vw_saldo_inscricao}\\
RF14 & \cod{tp_pagamento} e atributos opcionais & vencimento obrigatório para boleto e empenho\\
RF15 & Política de reembolso na Edição; situações de cancelamento & cancelamento em cascata nos itens e atividades; reembolso calculado pela política\\
RF16 & Vínculo Institucional com início e fim & vínculo com período\\
RF17 & Benefício, fraca de Instituição, e Facilita & vigência, instituição associada, vínculo vigente e limites verificados por trigger\\
RF18 & Trabalho, fraca de Edição, e Assina & pelo menos um autor por trabalho\\
RF19 & Cobre e Sobretaxa & cobertura só por autor em categoria de autor; excedente gera taxa\\
RF20 & \cod{dt_credenciamento} em Inscrição & credenciamento só em inscrição paga ou concluída\\
RF21 & \cod{dt_presenca} em Inscrição em Atividade & presença só em atividade confirmada\\
RF22 & Certificado com Certifica, Emite e Atesta & elegibilidade por tipo; um certificado válido por finalidade\\
RF23 & Termo com uma linha por versão; Responde & versão publicada não muda; resposta não é alterada\\
RF24 & Organiza com \cod{tp_papel} & chave (pessoa, edição, papel)\\
RF25 & Auditoria e Executa & trigger em inscrição, documento, pagamento e certificado\\
RF26 & Consultas sobre o modelo & doze visões de relatório (Seção~\ref{sec:relatorios})\\
RF27 & Valores congelados, preço e termo imutáveis & bloqueio de alteração em preço usado, termo, auditoria e políticas\\
RF28 & Chaves, domínios e regras entre linhas & PK, FK, UNIQUE, CHECK e triggers; 73 testes\\
RF29 & Inscrição em Atividade & bloqueio da linha da atividade durante a inscrição\\
RF30 & Poucos dados pessoais; documento só com caminho & papéis, permissões por coluna, visão mascarada e anonimização\\
\end{longtable}

\subsection{Questões levantadas pelo enunciado}

A seção 1.5 do enunciado lista questões que a equipe deveria descobrir durante o projeto. A Tabela~\ref{tab:questoes} resume a decisão tomada para cada uma.

\begin{longtable}{@{}L{5.6cm} L{8.7cm}@{}}
\caption{Respostas às questões do enunciado}\label{tab:questoes}\\
\toprule
Questão & Decisão\\
\midrule
\endfirsthead
\toprule
Questão & Decisão\\
\midrule
\endhead
\bottomrule
\endfoot
Inscrição iniciada em uma faixa e paga depois do fim dela & Vale a faixa vigente na data da inscrição. O preço fica reservado por \cod{nr_dias_reserva} dias; se a inscrição não for paga até \cod{dt_validade_preco}, ela expira.\\
Preço calculado na hora ou registrado & Registrado. Os itens financeiros guardam o valor no momento da cobrança.\\
Histórico quando a organização muda um preço & Preço já usado não pode ser alterado nem excluído. Uma mudança é um novo preço.\\
Inscrição com acesso a vários eventos co-localizados & Autorrelacionamento Agrega entre edições; o acesso é derivado por consulta.\\
Duas inscrições disputando a última vaga & A linha da atividade é bloqueada durante a inscrição e a ocupação é contada antes de aceitar.\\
Troca de categoria depois de enviar comprovante & Permitida enquanto a inscrição não está paga. A troca muda o preço, marca os itens antigos como SUBSTITUIDO e cria os novos.\\
Pagamentos já feitos quando a categoria muda & Como a troca só ocorre antes da quitação, os pagamentos confirmados continuam valendo e são descontados do novo saldo.\\
Pagamentos parciais ou múltiplos & A inscrição aceita vários pagamentos e passa a paga quando os confirmados cobrem o saldo.\\
Pagamento correspondente à obrigação certa & O pagamento pertence a uma inscrição, e o código da transação externa é único.\\
Benefícios com limites por ano, evento e instituição & O benefício pertence à instituição e tem limite por período e por edição. O limite por evento foi retirado (RE6).\\
Trabalho coberto por inscrição de autor & A cobertura só é aceita se a inscrição for de um autor, da mesma edição e em categoria de autor.\\
Trabalho com autores em categorias diferentes & Basta uma inscrição de autor em categoria de autor; cada trabalho tem no máximo uma inscrição que o cobre.\\
Autor responsável pelo pagamento cancela a inscrição & A cobertura é liberada e o trabalho volta a aparecer como sem cobertura.\\
Presença suficiente para cada certificado & Participação exige credenciamento, quando a edição pede, e frequência mínima. Atividade exige presença registrada. Autoria e organização exigem o vínculo correspondente.\\
Mudança de regra de certificação e certificados já emitidos & As políticas da edição ficam bloqueadas depois da abertura das inscrições, e um certificado emitido só pode ser revogado.\\
Atributos historizados & Valores dos itens, preços usados, versões de termo, tentativas de documento e de pagamento, e valores anterior e novo na auditoria.\\
Dados calculados e dados armazenados & Total devido, saldo, situação de associação, ocupação e aptidão a certificado são calculados por visões. Valores cobrados são armazenados.\\
Regras por CHECK, UNIQUE e FK & Domínios, intervalos e unicidade dentro da linha ou da tabela.\\
Regras por transação, trigger ou procedimento & Regras que comparam linhas ou tabelas, como sobreposição, capacidade, limites e transições de situação, ficam em triggers. Expiração e anonimização são procedimentos.\\
Consultas que pedem índices & Ficam para o Marco 2, que trata de otimização. No Marco 1 o banco tem só os índices criados pelas chaves primárias e restrições de unicidade.\\
Dados pessoais com mais controle & CPF, passaporte, e-mail e nascimento. Papéis sem acesso direto a eles usam uma visão mascarada.\\
Exclusão ou anonimização sem perder dados financeiros & A pessoa é anonimizada, e não excluída; inscrições, itens e pagamentos continuam íntegros.\\
\end{longtable}

\subsection{Sistematização entre entidades, relacionamentos e requisitos}

A Tabela~\ref{tab:sistematizacao} cruza cada entidade com os relacionamentos de que participa e os requisitos que atende. Os relacionamentos estão descritos na Tabela~\ref{tab:relacionamentos}.

\begin{longtable}{@{}L{3.3cm} L{5.9cm} L{4.8cm}@{}}
\caption{Entidades, relacionamentos e requisitos}\label{tab:sistematizacao}\\
\toprule
Entidade & Relacionamentos & Requisitos\\
\midrule
\endfirsthead
\toprule
Entidade & Relacionamentos & Requisitos\\
\midrule
\endhead
\bottomrule
\endfoot
Pessoa & Efetua, Vínculo Institucional, Associa-se, Analisa, Assina, Organiza, Certifica, Executa & RF3, RF4, RF8, RF16, RF18, RF22, RF24, RF25, RF30; RE8, RE13\\
Evento & Realiza & RF1\\
Edição & Realiza, Agrega, Oferece, Estabelece, Programa, Publica, Organiza, Emite & RF1, RF2, RF5, RF6, RF10, RF15, RF18, RF22, RF24, RF27; RE2, RE3, RE9\\
Categoria de Inscrição & Oferece, Preço, Restringe & RF5, RF8, RF10; RE17\\
Faixa de Inscrição & Estabelece, Preço & RF6\\
Preço & associativa entre Categoria e Faixa; Enquadra & RF6, RF7, RF27; RE1, RE5\\
Atividade & Programa, Inscrição em Atividade, Restringe, Atesta & RF10, RF11, RF22; RE17\\
Instituição & Vínculo Institucional, Recebe & RF16, RF17\\
Vínculo Institucional & associativa entre Pessoa e Instituição & RF16; RE14\\
Benefício & Recebe, Facilita & RF17; RE6\\
Período de Associação & Associa-se, Origina & RF4, RF7\\
Inscrição & Efetua, Enquadra, Facilita, Origina, Inscrição em Atividade, Compõe, Transaciona, Garante, Cobre, Responde & RF3, RF9, RF12, RF13, RF15, RF19, RF20, RF23; RE2, RE5, RE11\\
Inscrição em Atividade & associativa entre Inscrição e Atividade; Tarifa & RF10, RF11, RF21, RF29; RE4\\
Item Financeiro & Compõe, Tarifa, Sobretaxa & RF12, RF19, RF27; RE1, RE5\\
Pagamento & Transaciona & RF13, RF14, RF15; RE15\\
Documento de Comprovação & Garante, Analisa & RF8; RE12\\
Trabalho Científico & Publica, Cobre, Assina, Sobretaxa & RF18, RF19; RE9\\
Termo & Responde & RF23\\
Certificado & Certifica, Emite, Atesta & RF22; RE16\\
Auditoria & Executa & RF9, RF25\\
\end{longtable}
